\Needspace{34\baselineskip}
\subsection{算例数据地图：现在用的是哪一套参数？}
\label{app:data-map}
下表汇总各组算例的输入、参数与章节入口，便于跨章复算。
三维贯通主例沿用\hyperref[app:transformer-intermediates]{附录~\ref*{app:transformer-intermediates}} 的同一套初始参数，训练时从这套数值出发更新；
其他小例子各自给出数据，不与主例混用。
上下文对照中的“冬天”“窗”用于构造不同前文，“谢了”只作为二分类的输出类别，均不加入主线五项词表；二分类输入使用“春天／冬天／花”的三维 one-hot 编码。

\begin{center}\small
\renewcommand{\arraystretch}{1.3}
\tagpdfsetup{table/header-rows={1}}
\begin{tabularx}{\linewidth}{@{}>{\raggedright\arraybackslash}p{34mm}>{\raggedright\arraybackslash}p{45mm}>{\raggedright\arraybackslash}X@{}}
\toprule
算例／入口 & $X$、$W$ 或记录是什么 & 用途与对应章节\\\midrule
标量与两层网络 & $x,w,b$；再到 $xW_1,W_2$ & 函数、损失与反传；第~\ref*{intro:functions-and-loss}--\ref*{intro:mlp}~章\\
二维 Q/K/V & 手设向量与投影；各节写出 $X,W_Q,W_K,W_V$ & 逐项核对点积、读取与形状；第~\ref*{intro:attention-qkv}--\ref*{intro:self-attention}~章\\
四维双头 & 两个位置、每行四维；拆为两个二维头 & 拼接和输出投影；第~\ref*{intro:position-multihead}~章\\
三维贯通例 & $X=E_{[2,4,5]}+P_{\rm pos}$；参数见\hyperref[app:transformer-intermediates]{附录~\ref*{app:transformer-intermediates}} & 前向、联合训练、生成；第~\ref*{intro:language-model-task}、\ref*{intro:transformer-architecture}--\ref*{intro:real-models}~章\\
上下文与选择消融 & 局部 one-hot／相同末 Token；冻结或训练 Q/K & 区分读取前文与选择记录；第~\ref*{intro:context-gap}、\ref*{intro:language-model-training}~章\\
ICL 前驱与后继记录 & 四维记录投影为二维 $K,V$；从序列生成 & 示例改变、权重固定；第~\ref*{sec:icl-induction}~节\\
三分类奖励 & logits 三维；四个采样动作与固定优势 & 优势如何改变 logits；第~\ref*{concept:rlvr}~节\\
二维 RAG & $q=(1,0.5)$；三行文档矩阵 $D$ & 检索排序与证据覆盖；第~\ref*{intro:rag}~章\\
二维分数迷你版 & 两位置、两维；参数见\hyperref[app:rational-mini]{附录~\ref*{app:rational-mini}} & 20 分钟前向与反向；\hyperref[app:rational-mini]{附录~\ref*{app:rational-mini}}\\
联想记忆 & 各节给出 Key/Value，写入 $\mathcal M$ & 外积、串扰、Delta 与非线性；第~\ref*{intro:key-value-outer-product}--\ref*{intro:memory-capacity}~章\\
\bottomrule
\end{tabularx}
\end{center}
手算数据与脚本的对应入口见\hyperref[app:reproduction-scripts]{附录~\ref*{app:reproduction-scripts} 的复算脚本表（第~\pageref*{app:reproduction-scripts}~页）}。
